% !TEX root = ~/Code/TN_LOWE_RECONSTRUCTION/main.tex
% This is the frontmatter for the text.Ishould include the title, author, date, and an abstract
\title{Low-Energy Neutrino Event Reconstruction in DUNE: Clustering, Calibration, and Energy Reconstruction}

% Add affiliation if needed
\newcommand{\ciemat}{\small \textit{Centro de Investigaciones Energéticas, Medioambientales y Tecnológicas (CIEMAT), Madrid, Spain}}

\author{\textbf{Sergio Manthey Corchado}\\
    \ciemat\\
}

\date{\today}
\maketitle
\begin{abstract}
    The \ac{DUNE} aims to characterise the two most energetic components of the low--energy flux of solar neutrinos ($^{8}\text{B}$ \& $hep$), as well as neutrinos from other astrophysical sources such as nearby supernovae or the diffuse supernovae neutrino background. The sensitivity to the charged current channel, with orders of magnitude higher cross--section than the elastic scattering interaction, together with DUNE's high exposure and great energy resolution, present an exceptional advantage in understanding the sharp energy characteristics of this neutrino flux. This work presents a detailed study of low-energy charged-current neutrino event reconstruction and energy resolution performance in the presence of significant radiological and external backgrounds. The development and validation of a clustering algorithm that identifies primary ionisation signals and associates adjacent energy depositions to account for spatially correlated activity is described. A dedicated energy calibration strategy is employed to recover the true neutrino energy. The impact of various background scenarios on the reconstruction is assessed, and mitigation strategies are discussed. We report on the energy resolution achieved across a range of true neutrino energies, highlighting the algorithm’s effectiveness in maintaining reconstruction fidelity in the low-energy regime. These results demonstrate the robustness of the reconstruction approach and its potential contribution to the broader low--energy DUNE physics program.
\end{abstract}
\newpage
\tableofcontents
% \listoffigures
% \listoftables
\newpage